\section{Results and scope}
\label{intro4-section}

Stability relative to an interacting vacuum is a comparison of two
correlated ground energies. A lower bound proportional to spatial
volume does not supply that comparison. This paper proves such a
comparison for a compact Hamiltonian lattice gauge theory and for a
separately specified momentum-cutoff model. It then distinguishes
the resulting thermodynamic and low-energy consequences from
statements requiring new estimates.

\subsection{The gauge-invariant lattice result}

Fix $a>0$, $g^2>0$, $m\ge0$, and $r\ge0$. On the periodic cubic
lattice $X=(\Z/n\Z)^3$, take one four-component Wilson field and
compact $U(1)$ links with electric spectrum $\Z$. The full fermion
Fock space is used; there is no no-pair projection. The Hamiltonian,
including its raw fermion sea energy, is specified in
\eqref{lat-hamiltonian}. The physical space with integer static
source $\nu$ obeys every constraint
$\operatorname{div}E_x+n_x-2=\nu_x$ and is nonempty exactly when
$|\sum_x\nu_x|\le2n^3$.

The main source comparison is
\begin{equation}
 E_\nu\ge E_0+m\left|\sum_x\nu_x\right|.
 \label{intro4-main}
\end{equation}
Here $E_\nu$ includes the sources' electric self-energy, and $E_0$
is the exact interacting vacuum energy in the same box. No photon
mass, electric-flux truncation, weak-coupling assumption, or continuum
limit is used. For positive labeled nuclei of bounded charge,
the fixed per-species subtraction
\[
 S_\infty=\frac{g^2}{2a}\mathcal G_\infty(0)\sum_iZ_i^2
\]
gives
\begin{equation}
 E_\nu-S_\infty-E_0
 \ge mZ_{\rm tot}-\frac{\pi\sqrt3\,g^2}{16a}Z_{\max}^2K .
 \label{intro4-subtracted}
\end{equation}
This normalization is independent of state, positions, and volume.
A periodic-diagonal convention and its slightly different explicit
constant are also proved. Lattice sites resolve positions at scale
$a$; this is not an unregularized continuum point-source claim.

The mechanism is a positive centered finite-CAR determinant, not
positivity of arbitrary fermion matrix elements. The positive Wilson
mass makes the Dirac $\beta$ form decrease along the imaginary-time
propagator. The electric heat kernel and the local gauge projection
then turn each source into a unit-modulus character. A quantitative
Fourier bound proves the rest-energy strengthening in
\eqref{intro4-main}. The finite Fourier prefactor is retained before
the zero-temperature limit. The general transfer estimate is proved
in Lemmas~\ref{lem:mass-gap-annulus}--\ref{lem:mass-gap-fourier};
these forward references do not assume the lattice result.

\subsection{Thermodynamics and a sharply scoped low-energy result}

At fixed $a,g^2,m,r$, the unrestricted auxiliary ground energy equals
the physical vacuum energy. Cutting the auxiliary tensor-product
space therefore proves existence of the vacuum energy density,
with an explicit surface correction. No factorization of the
Gauss-constrained physical Hilbert space is assumed. For a fixed
admissible periodic source cell, a screening-flow trial state gives
two-sided source-relative bounds. A separate tiling argument gives
a thermodynamic limit for open boxes with zero boundary flux.
Identifying this density with the periodic-source limit still
requires a boundary-sector comparison; the missing statement is
specified in Section~\ref{therm4-section}.

At frozen links the positive Wilson bands admit an exact
Schur-complement comparison with
\[
 K(U)=R(U)+\frac{S(U)^2}{2m}.
\]
The addition-energy error is proportional to the number of added
particles, uniformly in lattice size and frozen links.
This does not integrate out pairs in the quantized gauge theory.
For the unimproved discretization, the small-momentum coefficients are
\[
 m_{\rm kin}=\frac{m}{1+rma},\qquad
 g_{\rm eff}=\frac{2}{1+rma}.
\]
The magnetic statement is a controlled smooth-field, low-momentum
matching coefficient, not a finite-coupling measurement of the
interacting anomalous moment. Explicit scale inequalities include
nuclear charge. A bounded magnetic clover modification can restore
the tree-level coefficient while preserving determinant positivity
in a stated bare-parameter range.

Two proposed dynamical shortcuts are false. Every eigenvector of
the compact model has infinite electric Fourier support, and no
vacuum eigenvector is the bare on-site sea. Thus a source-only bound
on bare high-flux tails or bare pair number cannot have zero
right side in the source-free vacuum. Smooth electric localization,
an explicit virtual-pair density bound, and the local-annihilation
remainder lemma are proved instead. The remaining dressed
pair-block resolvent and source-relative remainder estimates are
stated as open hypotheses, not consequences of locality.

\subsection{The separate momentum-cutoff comparison}

Sections~\ref{sec:model}--\ref{pos-section} retain continuum point
nuclei but use a sharp finite fermion momentum band, finite transverse
photon modes, and an explicit periodic Coulomb convention.
The harmonic photon zero mode is fixed. A positive effective mass
gives the canonical vacuum comparison, and heat-kernel electrostatics
gives $E(R,Z)-E^{\vac}\ge-C_\Lambda K$ in the stated box-size range.
The photon-coordinate limit is justified in quadratic-form and
trace-norm terms.

For the specified free-sea four-field Wick ordering, a sufficient
criterion uniform over all boxes permits $\Lambda_f/m<5168.21$ at
$\alpha=1/137$. The sharper logarithmic estimate has a mesh error;
an explicit coarse-box example prevents an all-volume exponential
window by that estimate. Positive density-product ordering is a
different Hamiltonian prescription with no such mass window.
Neither prescription is asserted to be a renormalized physical
QED construction.

The lattice and momentum models have different regulators and
different self-energy conventions. Their constants are not
interchanged. A charge gap is not a neutral pair gap; an extensive
vacuum density is not binding; and none of the present source costs
is uniform in the ultraviolet cutoff at fixed physical parameters.
